\maketitle
\head{本版编排与收录说明}
\noindent\textbf{这是一份按内容顺序整理的中文工作译稿，不是全部已上传材料的完结全译本。}第I、II章及第III章的两个已出版分册正文已经连续收入；第0章当前收入§§1--18及§19的19.0--19.6；第IV章主体、第0章19.7--19.10及§§20--23尚未连续译入；第V卷英文预备笔记只收入已译出的部分。未译内容没有用原文扫描页、提要或空白标题代替译文。

\section*{正文次序}
本版不再按上传时间、译稿批次或独立PDF页序排列，而按原著内容编号组织：导言；第0章预备知识；第I章概形的语言；第II章若干类态射的初等整体研究；第III章凝聚层的上同调研究；后续待译范围说明；独立编排的第V卷预备笔记已译部分；最后为各原分册的书末资料。第III章第一部分的§§1--5与第二部分的§§6--7直接衔接，不再被第一部分的索引隔开。

第0章原来分载于第一卷、第三卷第一部分、第四卷第一部分。本版将已译的§§1--7、§§8--13、§§14--18和§19已译部分归并在一起，但保留 $0_{\mathrm I}$、$0_{\mathrm{III}}$、$0_{\mathrm{IV}}$ 等出处记号。它们标示原出版分册，并非重复的三章。这样的安排依内容编号，而不声称每一节的全部前提均已在前文出现；原著跨章引用仍按原号保留。

\section*{实际收录范围}
\begin{longtable}{@{}p{.25\linewidth}p{.67\linewidth}@{}}
\toprule 材料 & 已收入的连续译文及状态\\\midrule\endhead
第一卷（1960） & 导言；第0章§§1--7；第I章§§1--10，止于（10.15.7）；原书书末资料。承接既有V2校订源码。\\
第二卷（1961） & 第II章§§1--8，止于（8.14.14）；原书参考文献、索引、目录、勘误清单一。其书末资料沿用上一合订版依据底本恢复的译排来源。\\
第三卷第一部分（1961） & 第0章§§8--13；第III章§§1--5，止于（5.5.2）；原书书末资料。\\
第三卷第二部分（1963） & 第III章§§6--7，止于（7.9.14），连同其证明、记号索引、术语索引、原目录及勘误清单二。第三版已补齐（6.7.1）至本分册末尾；本版保留。\\
第四卷第一部分（1964） & 第0章补篇章首说明、§§14--18及§19的19.0--19.6，止于（19.6.7）末尾。本版新增（19.5.1）--（19.6.7）；19.7以后尚未译入。\\
第V卷英文预备笔记 & \texttt{EGAV1-2.ps} 的六页可读材料；\texttt{EGAVaiello.pdf} 的卷首说明及“超平面截面与锥投影”§§1--6。余下内容不计入已译范围。\\
\bottomrule
\end{longtable}

\section*{译文、勘误与核查的关系}
正文采用本轮已经核实的正确陈述、公式和证明；原文错误写法、修改理由及依据记入脚注。作者正式勘误与本次译者补证分开，较长反例与验证收入独立核查报告。尚未确认的疑点不擅自改成确定结论；原著本来省略的证明、预备笔记中未写完或无法辨认的内容，不凭通常知识补造成“原文”。本轮新增注号以G7开头，旧注号保留以便追踪。

这一编辑原则已落实在本轮新增的19.5--19.6，并回溯修改了（19.3.5，v）、（19.3.10.1）、（19.3.12）和（19.4.1）。\textbf{它不表示此前全部旧译文中的历史核查意见均已逐项回写正文。}其他既有部分继续保存原稿；其校订状态须参看各条脚注和历史报告，后续仍需逐项核定。

对作者已撤销的（19.3.12），正文仅保留编号和撤销说明，原陈述及问题写入脚注；完整历史陈述与原证明保存在报告附录，不再占据有效论述的正文位置。作者勘误清单二、清单三相关选译仍在书末收录。未经核实的替代结论没有补造。

第V卷上传材料的标题明确标为预备笔记的翻译与编辑，英译编者为Piotr Blass与Joseph Blass。它们包含重组建议、未完成句子及待补证明，不按正式刊行的第I--IV章同等状态标示。与英文合并稿重合的PS载体不重复计作新译文。

\section*{第七版新增范围与回溯修订}
本版从第六版实际接续点（19.5.1）起，译入19.5“形式光滑性与相伴分次环”以及19.6“域上的代数的情形”，止于（19.6.7）末尾。对应第四卷第一部分分册印刷第90--104页（上传PDF第89--103页）。本轮包括一般定理的全部分段证明、Cohen定理、系数域和几何正则性判据；19.7章首及（19.7.1）仍为下一接续点，未算作译完。

（19.6.2）在正文补明相容的子域嵌入 $k\subset A$；原文仅给 $k\subset K$ 的问题及 $\mathbf Z_p$ 反例放在脚注和报告中。（19.5.6.7）的核比较步骤换成有限滤过逐级消去论证；（19.6.6）以Galois作用在局部分量上的传递性及忠实平坦下降修补原证明。后两项是证明修补，不宣称原命题结论有反例。

第0章其余部分、第I--III章及预备笔记的旧正文，除上列回溯项外没有进行一轮新的逐句复审。本次新核查范围只包括新增内容及上述直接复核的相关条目。累计核查报告中“第六版”“第五版”等章节是历史档案，其中旧的正文状态须以本版实际源码为准。

\section*{版式与使用}
采用 \texttt{\string\documentclass[UTF8,a4paper]\{ctexbook\}}，由文档类决定默认中文字体、字号、版心、标题与双面分页。不另设页边距，不用自定义页眉。为保留原著结构，仅保留原章号、节号与公式号；默认双面排版所需的空白偶数页不删除。前置部分使用罗马数字页码，其后连续编号；总目录指向本合订本页面，各原著目录的页码仍指原版。

PDF可用于对照研读；可独立编译的单文件TeX与分文件工程同时提供。旧稿内容的编排归并与技术校验，不等于本次将旧稿的每一句译文或每一个证明重新独立审定。累计核查报告单独交付，历史核查意见带有批次说明；对同一事项，本版直接核得的作者勘误和明确论证须与历史意见对照阅读，不能将历史“待核”标签当作最新状态。
\tableofcontents
